% ##############################################################################
\section{Authorization Drift}
\label{sec:problem}

Most SaaS platforms start with a clean idea. The backend enforces access. The frontend checks enough state to avoid showing nonsense. That works while the product is small.

Then tenants diverge. One customer buys Grid, another buys Horizon, a demo tenant gets DataMap, and support engineers need temporary visibility into each of them. A few flags become a few route guards. A layout component grows tenant-specific branches. A module entry appears in one sidebar but not another. The backend remains the final authority, yet the UI has become a second policy system without the same discipline.

We call the gap between these two models the \emph{representation gap}. It is the distance between what the backend will permit and what the UI presents as available.

\subsection{Requirements}
\label{sec:requirements}

A multi-tenant UI policy system must satisfy five requirements.

\begin{description}[leftmargin=2.2em, style=nextline]
  \item[R1. No stale gates.] Policy-gated UI elements shown to a user must correspond to backend operations the system will actually permit for the user's current tenant and role.
  \item[R2. Configuration over redeploy.] Enabling a tenant or changing a module entitlement should be a reviewed data change, not a code release.
  \item[R3. Safe impersonation.] Support engineers need the tenant's perspective without shared credentials and without bypassing the authorization model.
  \item[R4. Defense in depth.] The UI must never be the only barrier. Backend APIs authorize requests even when a browser is tampered with.
  \item[R5. Support at scale.] Operators need scoped concurrent sessions and searchable audit trails, especially when several engineers are debugging different tenants at the same time.
\end{description}

Traditional RBAC handles a useful slice of this problem \cite{sandhu_rbac}. ABAC gives a richer model for attribute-derived decisions \cite{nist_abac}. Feature flags help with release control \cite{fowler_feature_toggles}. Policy engines such as OPA, Cedar, and OpenFGA centralize authorization decisions \cite{opa_docs,cedar_docs,aws_verified_permissions,openfga_docs}.

None of those choices automatically tells the frontend what world to render. They answer what is allowed. They do not prescribe how permission decisions become navigation, layout, module configuration, and support workflow. That missing layer is where \uiPolicy{} sits.

\begin{table}[H]
  \centering
  \small
  \begin{tabularx}{\textwidth}{X c c c}
    \toprule
    \textbf{Need} & \textbf{RBAC / ABAC} & \textbf{Feature flags} & \textbf{\uiPolicy{}} \\
    \midrule
    Backend authorization & \cmark & \xmark & \cmark \\
    Tenant-level variability & Partial & \cmark & \cmark \\
    UI derived from policy & \xmark & Partial & \cmark \\
    Config-only tenant changes & Partial & \cmark & \cmark \\
    Auditable view-as support & \xmark & \xmark & \cmark \\
    Defense against stale UI gates & \xmark & \xmark & \cmark \\
    \bottomrule
  \end{tabularx}
  \caption{Where \uiPolicy{} fits. It does not replace authorization frameworks. It uses their decisions to make the rendered interface truthful.}
  \label{tab:comparison}
\end{table}

\subsection{The Truthful UX invariant}
\label{sec:truthful_ux}

For policy-gated surfaces, we use the following invariant.

\begin{quote}
\textbf{A rendered UI affordance must be backed by a server-side policy decision in the same effective tenant and role context.}
\end{quote}

The invariant is intentionally narrow. It does not claim every pixel on the page is authorization-sensitive. A help link or account menu may be universal. It does claim that navigation, modules, layouts, and policy-gated actions should not be rendered from a parallel frontend guess.
